% Préambule commun aux cours IPR Scholar.
\documentclass[11pt,a4paper,oneside]{article}
\usepackage[T1]{fontenc}
\usepackage[utf8]{inputenc}
\usepackage[french]{babel}
\usepackage{lmodern}
\usepackage{microtype}
\usepackage[a4paper,margin=2.35cm]{geometry}
\usepackage{amsmath,amssymb,mathtools}
\usepackage{booktabs,longtable,array}
\usepackage{xcolor}
\usepackage{hyperref}
\usepackage{tikz}
\usetikzlibrary{arrows.meta,positioning,shapes.geometric,calc}
\hypersetup{colorlinks=true,linkcolor=blue!55!black,citecolor=blue!55!black,urlcolor=blue!55!black}
\definecolor{iprblue}{HTML}{1F4E79}
\definecolor{iprgreen}{HTML}{2D6A4F}
\newcommand{\IPR}{\textsc{ipr}}
\newcommand{\coursmeta}[2]{\author{Projet \IPR{} -- documentation Scholar}\date{#1}\maketitle\begin{center}\small #2\end{center}\vspace{0.5cm}\tableofcontents\clearpage}
\newcommand{\resultat}[1]{\paragraph{À retenir.} #1}

\title{Cours 2 -- Exposition, contraste et masques de tonalité}
\begin{document}
\coursmeta{3 octobre 2026}{Approches photométriques utilisées par les réglages Exposition, Contraste, Ombres, Noirs, Hautes lumières et Blancs.}

\section{Théorie}
Une valeur d'exposition $E$ est logarithmique : un incrément d'un IL/EV double l'énergie.
\begin{equation}
\mathbf{p}'=2^E\mathbf{p}.
\end{equation}
Le contraste ne doit pas être une simple multiplication autour de zéro. En photographie, le gris moyen est un pivot perceptuel usuel, voisin de $18\%$ en lumière linéaire. \IPR{} emploie une courbe en S asymétrique qui conserve ce pivot : sous le pivot, une loi de puissance agit sur le \emph{toe}; au-dessus, la courbe utilise une épaule indépendante. Pour une luminance $Y\le p$, une forme représentative est $Y'=p(Y/p)^a$.

Les réglages de plage tonale utilisent une fenêtre douce, plutôt qu'un seuil. Cette logique relève des transformations ponctuelles et des fonctions de transfert étudiées en traitement d'image~\cite{gonzalezWoods}. La fonction utilisée est
\begin{equation}
S(e_0,e_1,x)=t^2(3-2t),\quad t=\operatorname{clamp}\left(\frac{x-e_0}{e_1-e_0},0,1\right).
\end{equation}
Elle est continue avec pente nulle aux bornes : les transitions entre ombres et tons moyens ne créent pas de bandes.

\section{Application dans \IPR{}}
Chaque correction calcule d'abord $Y$, puis modifie une luminance cible $Y'$ avant de remettre l'image à l'échelle : $\mathbf{p}'=\mathbf{p}\,Y'/Y$ si $Y>10^{-6}$. Les noirs sont confinés au pied de courbe; les ombres couvrent une plage plus large; les blancs utilisent aussi le maximum RGB afin de réagir aux canaux proches de la saturation. Les masques sont évalués sur la luminance d'entrée afin qu'un réglage ne déplace pas lui-même sa zone d'action.

\section{Exemple applicatif}
Considérons un pixel neutre $Y=0{,}04$. Avec \emph{Ombres} positif, le masque des ombres est proche de un et la correction le remonte sans éclaircir une zone à $Y=0{,}80$. Avec \emph{Noirs} positif, une forme $\sqrt{u}(1-u)$, où $u=Y/0{,}14$, laisse le noir absolu à zéro : l'absence de lumière ne devient pas un voile gris.

\section{Interprétation des constantes}
Le pivot $0{,}18$ a une justification photographique. Les bornes des masques ($0{,}10$, $0{,}56$, $0{,}70$, etc.) et les amplitudes sont des paramètres de rendu calibrés sur le corpus de validation Lightroom/IPR : elles ne sont pas des constantes universelles. Elles doivent être renommées et regroupées lorsqu'une calibration automatisée sera disponible.

\resultat{Exposition est un gain physique; tonalité et contraste sont des choix de \emph{tone mapping} continus, évalués par comparaison visuelle et métrique.}
\bibliographystyle{plain}\bibliography{references}
\end{document}
